\documentclass[a4paper,11pt]{jsarticle} % 日本語環境の場合
% 英語環境の場合は \documentclass[a4paper,11pt]{article} とし、日本語パッケージを外してください

\usepackage{amsmath}
\usepackage{amssymb}
\usepackage{bm}
\usepackage[dvipdfmx]{graphicx}
\usepackage{booktabs}
\usepackage{url}

\title{ドローン周辺流場における非対称性の定量評価手法}
\author{}
\date{}

\begin{document}

\maketitle

\section{目的}
壁面効果（地面効果および側壁効果）によって生じるドローン周辺の空気流れの非対称性を数値化し、機体の安定性への影響を定量的に評価することを目的とする。

\section{評価指標の定義}

\subsection{領域平均流速による非対称指数 (Asymmetry Index)}
プロペラ下方の特定断面において、壁側（Wall side）と開放側（Open side）の平均流速を比較し、その偏りを指数化する。

\begin{equation}
  I_{asym} = \frac{|\bar{V}_{Wall} - \bar{V}_{Open}|}{\bar{V}_{Total}}
\end{equation}

\begin{itemize}
  \item $I_{asym}$: 非対称指数（0で完全対称、値が大きいほど非対称）
  \item $\bar{V}_{Wall}$: 壁側領域の平均垂直流速 [m/s]
  \item $\bar{V}_{Open}$: 開放側領域の平均垂直流速 [m/s]
  \item $\bar{V}_{Total}$: 全領域の平均垂直流速 [m/s]
\end{itemize}

\subsection{流束重心偏位 (Displacement of Flow Center)}
プロペラから吹き下ろされる風の「重心位置（圧力中心に近い概念）」が、幾何学的中心からどの程度ズレているかを評価する。

\begin{equation}
  X_{center} = \frac{\int (x \cdot w) dA}{\int w dA}
\end{equation}

\begin{itemize}
  \item $x$: 機体中心からの水平距離座標
  \item $w$: 垂直方向流速成分 ($Velocity\_w$)
  \item $dA$: 微小面積
\end{itemize}

\subsection{流れの偏角 (Flow Deflection Angle)}
壁の存在により、本来垂直であるべき吹き下ろし風が何度傾いたかを算出する。

\begin{equation}
  \theta_{def} = \arctan \left( \frac{\bar{V}_{Horizontal}}{\bar{V}_{Vertical}} \right)
\end{equation}

\begin{itemize}
  \item $\bar{V}_{Horizontal}$: 水平方向の平均流速成分
  \item $\bar{V}_{Vertical}$: 垂直方向の平均流速成分
\end{itemize}
